\section{Step 49 Compare with an independent finite reference}\label{h263-transform-s09}
Run public feedback against your current source:
\begin{verbatim}
python3 construction/h263-transform-construction/grade.py \
    my-transform.asm --cli /path/to/mars-cli.jar \
    --report transform-results.json
\end{verbatim}
The 96 cases check both directions, modes and byte orders, mixed and extreme
patterns, exact zeros, padded/wide strides and rejected metadata or sample domains.
The reference computes floating-point matrix products independently of the
provided fixed-point tables. It checks every output, source access, saved state
and guard, accepting an accurate alternative precision or private frame.
Reports require a new path.

Try both flat and worked ordinary callers. A single printed first coefficient
is useful for debugging but does not test sixty-four positions. At the return
breakpoint, inspect the whole signed halfword block and the count before the
caller replaces return registers with syscall arguments.

\paragraph{Observe} Finite checks accept both supplied accurate precisions.
They reject wrong signed products, narrow second products, incorrect scaling,
wrong intermediate placement and omitted sample-domain checks.
\Needspace{8\baselineskip}
\paragraph{Check} Why use an independent reference instead of only a round trip?
\begin{itemize}
  \item[A.] A round trip guarantees every coefficient convention is correct.
  \item[B.] The reference must copy the student's intermediate registers.
  \item[C.] A forward and inverse pair can share the same normalization error.
\end{itemize}
